\unitopen{16}
\begin{candos}
\item zdecydować, czy użyć \de{haben} czy \de{sein} w Perfekcie;
\item znać imiesłowy najczęstszych czasowników nieregularnych (\de{gegangen, gefahren, gekommen, gegessen, geschrieben, \dots});
\item opowiedzieć o podróży lub urlopie w kilku zdaniach;
\item zapytać \de{Wohin bist du gefahren?} i odpowiedzieć.
\end{candos}

\section{Słownictwo}
\vocabtable{16}

\section{\de{sein} czy \de{haben}?}
\de{sein} tworzy Perfekt czasowników \emph{ruchu z miejsca na miejsce} i \emph{zmiany stanu}, oraz \de{sein}, \de{werden},
\de{bleiben}: \de{Ich bin nach Zürich gefahren.} \de{Er ist eingeschlafen.} \de{Wir sind zu Hause geblieben.} \de{Es ist kalt geworden.}
Wszystkie pozostałe używają \de{haben}: \de{Ich habe gearbeitet.} \de{Ich habe einen Film gesehen.}

\begin{center}\small
\begin{tabular}{@{}lll@{}}
\toprule
\textsf{\textbf{Bezokolicznik}} & \textsf{\textbf{Partizip II}} & \textsf{\textbf{Pomocniczy}}\\
\midrule
\de{gehen} & \de{gegangen} & \de{sein}\\
\de{fahren} & \de{gefahren} & \de{sein} (z biernikiem pojazdu: \de{haben}: \de{Ich habe das Auto gefahren})\\
\de{kommen} & \de{gekommen} & \de{sein}\\
\de{fliegen} & \de{geflogen} & \de{sein}\\
\de{bleiben} & \de{geblieben} & \de{sein}\\
\de{sein} & \de{gewesen} & \de{sein}\\
\de{essen} & \de{gegessen} & \de{haben}\\
\de{trinken} & \de{getrunken} & \de{haben}\\
\de{sehen} & \de{gesehen} & \de{haben}\\
\de{lesen} & \de{gelesen} & \de{haben}\\
\de{schreiben} & \de{geschrieben} & \de{haben}\\
\de{sprechen} & \de{gesprochen} & \de{haben}\\
\de{nehmen} & \de{genommen} & \de{haben}\\
\de{geben} & \de{gegeben} & \de{haben}\\
\de{finden} & \de{gefunden} & \de{haben}\\
\de{verstehen} & \de{verstanden} & \de{haben} (bez \de{ge-})\\
\bottomrule
\end{tabular}
\end{center}
Czasowniki nieregularne \emph{musisz} znać w formie bezokolicznik--imiesłów. Ucz ich w parach (Anki: karta
\de{gehen, ging, ist gegangen} jako jedna). Kończą się na \de{-en} (często ze zmianą samogłoski).

\begin{falle}
Polski „byłem w Zurychu” to \de{Ich \textbf{war} in Zürich} (Präteritum \de{sein}, jednostka~17) lub \de{Ich bin in Zürich
\textbf{gewesen}}; „pojechałem do Zurychu” to \de{Ich \textbf{bin} nach Zürich \textbf{gefahren}}. Nie używaj \de{haben}
z czasownikami ruchu: \de{Ich habe nach Zürich gefahren} jest błędem. Pułapka drugiego stopnia: \de{schwimmen}
bierze \de{sein}, gdy płyniesz dokądś (\de{Ich bin ans Ufer geschwommen}), a \de{haben}, gdy czynność trwa (\de{Ich habe eine Stunde geschwommen}).
\end{falle}

\begin{szwajcaria}
Wakacje Szwajcarów to \de{Ferien} (nie \de{Urlaub}, chociaż \de{Urlaub} też zrozumieją): \de{Wir sind in den Ferien
in die Berge gefahren.} Pytanie o urlop w pracy: \de{Wann hast du Ferien?} W CH ustawowe minimum to cztery tygodnie urlopu rocznie
\verify{(sprawdź w OR i w umowie)}.
\end{szwajcaria}

\begin{praca}
Konferencja lub delegacja: \de{Ich bin gestern nach Zürich geflogen und habe den Kunden getroffen.} \de{Wir sind
mit dem Zug gekommen.} W wiadomościach o błędach: \de{Der Build ist gestern fehlgeschlagen} (\de{fehlschlagen}: \de{sein}).
\end{praca}

\section{Dialogi}
\textbf{Dialog 1 --- po urlopie}
\begin{dialog}
\turn{A}{Wo warst du im Urlaub?}{Gdzie byłeś na urlopie?}
\turn{B}{Ich bin mit meiner Freundin nach Italien gefahren. Wir sind zehn Tage geblieben.}{Pojechałem z dziewczyną do Włoch. Zostaliśmy dziesięć dni.}
\turn{A}{Was habt ihr gemacht?}{Co robiliście?}
\turn{B}{Wir sind viel gelaufen, haben gut gegessen und sind im Meer geschwommen.}{Dużo chodziliśmy, dobrze jedliśmy i pływaliśmy w morzu.}
\turn{A}{Und wie war das Wetter?}{A jaka była pogoda?}
\turn{B}{Super! Aber ich habe meinen Koffer verloren.}{Super! Ale zgubiłem walizkę.}
\end{dialog}

\section{Ćwiczenia}
\begin{cwiczenia}
\exitem{Perfekt, \de{ich / nach Zürich / fahren}.}{\de{Ich bin nach Zürich gefahren.}}
\exitem{Perfekt, \de{wir / zu Hause / bleiben}.}{\de{Wir sind zu Hause geblieben.}}
\exitem{Perfekt, \de{er / ein Buch / lesen}.}{\de{Er hat ein Buch gelesen.}}
\exitem{Perfekt, \de{sie / Kaffee / trinken}.}{\de{Sie hat Kaffee getrunken.}}
\exitem{Perfekt, \de{ich / gestern / Deutsch / sprechen}.}{\de{Ich habe gestern Deutsch gesprochen.}}
\exitem{Perfekt, \de{wir / ins Kino / gehen}.}{\de{Wir sind ins Kino gegangen.}}
\exitem{Perfekt, \de{du / die Prüfung / bekommen} (dostać wynik).}{\de{Du hast die Prüfung bekommen.}} 
\exitem{Podaj \de{Partizip II}: \de{nehmen}.}{\de{genommen}}
\exitem{Podaj \de{Partizip II}: \de{finden}.}{\de{gefunden}}
\exitem{Podaj \de{Partizip II}: \de{sein}.}{\de{gewesen}}
\exitem{Popraw: \de{Ich habe nach Zürich gefahren.}}{\de{Ich bin nach Zürich gefahren.}}
\exitem{Przetłumacz: „Wczoraj poleciałem do Zurychu.”}{\de{Gestern bin ich nach Zürich geflogen.}}
\exitem{Zapytaj: „Gdzie byłeś na urlopie?”}{\de{Wo warst du im Urlaub?}}
\end{cwiczenia}

\section{Trening słownictwa (generowany)}
\drill{16}

\weekend{Opisz po niemiecku najlepszą podróż w swoim życiu w dziesięciu zdaniach, używając co najmniej
czterech czasowników ze \de{sein} i czterech z \de{haben} w Perfekcie. Nagraj i sprawdź pomocnicze.}
